\hypertarget{paiement-et-gestion-de-fraude}{%
\subsubsection{Paiement et Gestion de
fraude:}\label{paiement-et-gestion-de-fraude}}

\hypertarget{us1---cruxe9ation-dune-demande-de-paiement-par-qr-code}{%
\paragraph{US1 - Création d'une demande de paiement par QR
code:}\label{us1---cruxe9ation-dune-demande-de-paiement-par-qr-code}}

En tant que titulaire, je veux pouvoir créér une demande de QR code afin
de facilité la réception d'un paiements

\textbf{Critères d'acceptation :}

\begin{itemize}
\tightlist
\item
  l'application doit permettre a l'utilisateur de générer un QR code.
\item
  Le QR code doit contenir:

  \begin{itemize}
  \tightlist
  \item
    le montant à payer;
  \item
    le numéro du compte béneficiaire (IBAN);
  \item
    un message décrivant le paiement;
  \item
    la date et l'heure de création.
  \end{itemize}
\item
  L'utilisateur doit pouvoir téléchargé le QR code sous forme d'image
\item
  Le QR code doit etre valide pendant une durée maximal d'une heure
\item
  Le QR code devient automatiquement invalide:

  \begin{itemize}
  \tightlist
  \item
    après son expiration ;
  \item
    après sa première utilisation.
  \end{itemize}
\item
  Une fois le QR code validé par un payeur :

  \begin{itemize}
  \tightlist
  \item
    celui-ci choisit le compte source à débiter;
  \item
    le montant est transféré automatiquement vers le compte bénéficiare
    encodé dans le QR code.
  \end{itemize}
\end{itemize}

\hypertarget{us2---paiement-par-qr-code}{%
\paragraph{US2 - Paiement par QR
code:}\label{us2---paiement-par-qr-code}}

En tant que titulaire, je veux pouvoir payer par QR code afin de
facilité les paiements.

\textbf{Critères d'acceptation :}

\begin{itemize}
\tightlist
\item
  Le paiement ne doit pas être accepté si le QR code est invalide
  (expiré ou déjà utilisé).
\item
  Après une transaction réussie :

  \begin{itemize}
  \tightlist
  \item
    le solde du compte source doit être débité ;
  \item
    le solde du compte destinataire doit être crédité ;
  \item
    le QR code devient automatiquement invalide (usage unique).
  \end{itemize}
\end{itemize}

\hypertarget{us3---paiement-avec-ou-sans-contact}{%
\paragraph{US3 - Paiement avec ou sans contact
:}\label{us3---paiement-avec-ou-sans-contact}}

En tant titulaire, je veux pouvoir effectuer des paiements avec ou sans
contact afin de réaliser des achats de manière rapide et sécurisée.

\textbf{Critères d'acceptation :}

\begin{itemize}
\tightlist
\item
  L'utilisateur peut sélectionner un paiement parmi ceux disponibles
  dans l'application.
\item
  L'utilisateur peut choisir entre :

  \begin{itemize}
  \tightlist
  \item
    un paiement avec contact;
  \item
    un paiement sans contact.
  \end{itemize}
\item
  Pour un paiement avec contact :

  \begin{itemize}
  \tightlist
  \item
    la confirmation par PIN est obligatoire.
  \end{itemize}
\item
  Pour un paiement sans contact:

  \begin{itemize}
  \tightlist
  \item
    Aucun PIN n'est demandé tant que aucun suspicion de fraud n'est
    détecté.
  \end{itemize}
\item
  L'utilisateur doit pouvoir sélectionner le compte source à débiter
\item
  Le paiement doit etre refusé si le solde est insuffisant.
\item
  Le montant du paiement et le destinataire sont automatiquement
  récupérer depuis le paiement sélectionné
\item
  Apres une transaction reussit:

  \begin{itemize}
  \tightlist
  \item
    Le montant doit etre débité du solde du compte source;
  \item
    Le montant doit doit crédité au solde du compte destinataire;
  \item
    Les soldes sont immédiatemment mis à jour.
  \end{itemize}
\item
  L'institution financière doit pouvoir configurer :

  \begin{itemize}
  \tightlist
  \item
    Une limite par transaction;
  \item
    Une limite journalière;
  \item
    Une limite hebdomadaire;
  \item
    Une limite mensuelle.
  \end{itemize}
\item
  Le paiement est refusé lorsque l'une des limites configurées est
  dépassée
\end{itemize}

\hypertarget{us4--analyse-des-transactions}{%
\paragraph{US4- Analyse des
transactions:}\label{us4--analyse-des-transactions}}

En tant que System (Automatique), je veut que les transactions des
clients soit analyse afin que toute activite suspecte sur le compte
client soit detecter.

\textbf{Critères d'acceptation :}

\begin{itemize}
\item
  L'application doit enregistrer les métadonnées suivantes pour chaque
  transaction :

  \begin{itemize}
  \tightlist
  \item
    Montant, date, heure, mode de paiement, lieu de paiement, pays de
    paiement, destinataire.
  \end{itemize}
\item
  L'application doit conserver l'historique des transactions pour
  affiner la détection.
\item
  Le système doit identifier les anomalies suivantes :

  \begin{itemize}
  \tightlist
  \item
    Plusieurs paiements successifs dans un court intervalle de temps.
  \item
    Paiements effectués dans plusieurs pays différents dans un délai
    anormalement court.
  \item
    Montants inhabituels par rapport aux habitudes du client.
  \item
    Paiements vers des destinataires présents sur une liste noire.
  \end{itemize}
\item
  Le système doit tenir compte des suspicions précédentes pour ajuster
  son niveau d'alerte.
\item
  Lorsqu'un suspicion est détectée, l'utilisateur doit recevoir une
  notification.
\end{itemize}

\hypertarget{us5--gestion-de-fraude}{%
\paragraph{US5- Gestion de fraude :}\label{us5--gestion-de-fraude}}

En tant que Systeme (Automatique), je veux réagir automatiquement
lorsqu'une fraude est suspectée afin de protéger le compte client.

\textbf{Critères d'acceptation :}

\begin{itemize}
\tightlist
\item
  Le système doit attribuer un niveau de risque:

  \begin{itemize}
  \tightlist
  \item
    faible;
  \item
    moyen;
  \item
    fort.
  \end{itemize}
\item
  En cas de faible suspicion:

  \begin{itemize}
  \tightlist
  \item
    toute nouvelle transaction doit exiger une authentification (code
    PIN ou mots de passe).
  \item
    le système doit notifier le client
  \end{itemize}
\item
  En cas de moyen suspicion:

  \begin{itemize}
  \tightlist
  \item
    la transaction est bloquée;
  \item
    une notification est envoyée au client;
  \item
    le client peut confirmer ou refuser l'opération.
  \end{itemize}
\item
  En cas de fort:

  \begin{itemize}
  \tightlist
  \item
    le compte est temporairement bloqué;
  \item
    le client est contacté via un canal sécurisé.
  \end{itemize}
\item
  Si le client signal une fausse alerte

  \begin{itemize}
  \tightlist
  \item
    toute les restrictions appliquées sont supprimées
  \end{itemize}
\item
  Toute desicion prise doivent etre enregistrées dans l'historique des
  fraudes.
\item
  Chaque suspicion doit etre enregistrée dans un historique dédié.
\end{itemize}

\hypertarget{us6---consultation-de-lhistorique-des-transactions}{%
\paragraph{US6 - Consultation de l'historique des transactions
:}\label{us6---consultation-de-lhistorique-des-transactions}}

En tant que titulaire, je veut consulter mon historique de transaction
afin de consulté mes opération et détecté toute activité suspecte.

\textbf{Critères d'acceptation :}

\begin{itemize}
\tightlist
\item
  Toutes les transaction passé doivent figurée dans l'historique des
  transaction
\item
  L'application doit enregistrer pour chaque transaction:

  \begin{itemize}
  \tightlist
  \item
    date et l'heure;
  \item
    montant;
  \item
    le compte source;
  \item
    le compte destinataire;
  \item
    lieu de la transaction;
  \item
    le mode de paiement;
  \item
    statut du paiement;
  \end{itemize}
\item
  L'utilisateur doit pouvoir consulter cette historique a tout moments.
\item
  Chaque transaction doit afficher l'un des statuts suivant:

  \begin{itemize}
  \tightlist
  \item
    Reussie;
  \item
    Refusée;
  \item
    Bloquée;
  \end{itemize}
\item
  Toute les informations nécessaire a la détection de fraude doivent
  figuré et etre conservé dans l'historique;
\item
  Les transaction peuvent etre triées par date.
\item
  L'utilisateur peut consulter l'historique des transactions suspectes
  ainsi que leur niveau de risque.
\end{itemize}

\hypertarget{us7---signalement-des-faux-positifs-et-consultation-des-alertes}{%
\paragraph{US7 - Signalement des faux positifs et consultation des
alertes:}\label{us7---signalement-des-faux-positifs-et-consultation-des-alertes}}

En tant que titulaire, je veux pouvoir consulter les alertes de fraudes
détectées sur mon compte afin de comprendres les actions prise par le
systemes et déclarer les faux positifs.

\textbf{Critères d'acceptation :}

\begin{itemize}
\tightlist
\item
  Toute les transaction détécté comme suspecte de fraude doivent figurée
  dans l'historique.
\item
  Pour chaque alerte, les informations suivant doivent etre disponibles:

  \begin{itemize}
  \tightlist
  \item
    date et heure;
  \item
    transaction concernée;
  \item
    niveau de risque (faible,moyen et fort);
  \item
    action appliqué par le système.
  \end{itemize}
\item
  L'utilisateur doit pouvoir consulter l'historique des alertes de
  fraude.
\item
  L'utilisateur doit pouvoir indique qu'une alerte est un faux possitif.
\item
  Le système doit conserver les désicions de l'utilisateur afin
  d'améliorer les futurs controles.
\item
  L'etat de chaque alerte doit etre affiché:

  \begin{itemize}
  \tightlist
  \item
    en attente;
  \item
    confirmée comme fraude;
  \item
    faux possitif.
  \end{itemize}
\end{itemize}
